\RequirePackage{silence}
\WarningsOff[hyperref]
\input{preamble}

\renewcommand{\JournalYear}{2021}
\vsgtu{xxxx} %registration number
\FirstPage{1}
\LastPage{x}

\pdfcrop
%\draft

\begin{document}

\udc{XXX.XXX}
\msc{YYYYY, ZZZZZ, DDDDD}

\rutitle{Основное название статьи}
\rutitleshort{Короткое название статьи, если необходимо}

\entitle{The  title of the article}
\entitleshort{The short title of this article if necessary}

\ruauthor{А.~А.~Иванов\affil{1,2}, В.~В.~Смирнов\affil{2}}{И\,в\,а\,н\,о\,в~А.~В., С\,м\,и\,р\,н\,о\,в~В.~В.}
\enauthor{A.~A.~Ivanov\affil{1,2}, V.~V.~Smirnov\affil{2}}{I\,v\,a\,n\,o\,v~A.~V., S\,m\,i\,r\,n\,o\,v~V.~V.}

\ruaffil{Московский государственный университет имени М.~В.~Ломоносова,\\ 
Научно-исследовательский институт механики, \\
Россия, 119192, Москва, Мичуринский проспект, 1.
\and 
Самарский государственный технический университет,\\
Россия, 443100, Самара, ул. Молодогвардейская, 244.}

\enaffil{Lomonosov Moscow State University,\\ 
Institute of Mechanics, \\
1, Michurinsky prospekt, Moscow, 119192, Russian Federation.
\and 
Samara State Technical University,\\
244, Molodogvardeyskaya st., Samara, 443100, Russian Federation.}

\ruauthorsdetails{2}{% Количество авторов
\fullauthor{Андрей Алексеевич Иванов}
\CAuthor % Автор-корреспондент
\orcid{0000-0000-0000-xxxx} % ORCID ID автора 
\regalia{кандидат физико-математических наук, доцент} 
\position{старший научный сотрудник} 
\dept{лаб. упругости и пластичности} 
\email{anton@ya-mail.ru}
\fullauthor{Владимир Викторович Смирнов}
\orcid{0000-0000-0000-xxxx} % ORCID ID автора 
\regalia{доктор физико-математических наук, профессор} 
\position{заведующий кафедрой} 
\dept{каф. упругости и пластичности} 
\email{vlad_smirnov@ya-mail.ru}
}

\enauthorsdetails{2}{
\fullauthor{Andrew A. Ivanov}
\CAuthor % Сorresponding Author
\orcid{0000-0000-0000-xxxx} % ORCID ID
\regalia{Cand. Phys. \& Math. Sci., Associate Professor} 
\position{Senior Researcher} 
\dept{Lab. of Elasticity and Plasticity}
\email{anton@ya-mail.ru}
\fullauthor{Vladimir V. Smirnov}
\orcid{0000-0000-0000-xxxx} % ORCID ID avtora 
\regalia{Dr. Phys. \& Math. Sci., Professor} 
\position{Head of Department} 
\dept{Dept. of Elasticity and Plasticity}
\email[br]{vlad_smirnov@ya-mail.ru}
}

\rukeywords{ключевое слово, ключевая фраза, новая мысль}
\enkeywords{keyword, key phrase, new thought}

\ruabstract{Аннотация на русском языке должна содержать 9–15 предложений (не менее 200–250 слов). Разрешается структурировать аннотацию: объект исследования, цель, использованные методы и подходы, основные результаты. Аннотация должна быть информативной, отражать основное содержание статьи и результаты исследований, следовать логике описания результатов в статье и не содержать общих слов. Аннотация статьи на английском языке для иностранных учёных и специалистов является основным источником информации о содержании статьи и изложенных в ней результатах исследований. В этом случае аннотация выполняет функцию независимого от статьи источника информации. Поэтому к~её написанию необходимо относиться с должным вниманием. Аннотация на английском языке должна быть оригинальной (не быть калькой русскоязычной аннотации с дословным переводом), написанной качественным английским языком. Аннотация не должна содержать ссылок на литературу и аббревиатуры (если это возможно).}

\enabstract{An abstract is a succinct summary of a longer piece of work, usually academic in nature, which is published in isolation from the main text and should therefore stand on its own and be understandable without reference to the longer piece. It should report the latter's essential facts, and should not exaggerate or contain material that is not there. Its purpose is to act as a reference tool (for example in a library abstracting service), enabling the reader to decide whether or not to read the full text. Abstracts should contain no more than 250 words. Write concisely and clearly. 

The abstract should reflect only what appears in the original paper. The Abstract should:
\begin{itemize}
\item Describe the main objective(s) of the study;
\item Explain how the study was done, including any model organisms used, without methodological detail;
\item Summarize the most important results and their significance.
\end{itemize}
Abstracts should not include citations and abbreviations, if possible.}

%\date{\today}
%\revisiondate{\today}
%\accepteddate{\today}

%\ForPeerReview

\makerutitle

\newpage

\Section[n]{Введение}
Это шаблон статьи для представления рукописей статей в наш журнал.

Представляемая в журнал рукопись статьи должна быть законченным научным исследованием, нигде ранее не публиковавшимся и не представленным к публикации в других изданиях.

Рукопись статьи должна содержать новые научные результаты по приоритетным направлениям Самарского государственного технического университета, таким как «Дифференциальные уравнения и математическая физика», «Механика деформируемого твёрдого тела», «Математическое моделирование, численные методы и комплексы программ». 

Этот шаблон был разработан в сотрудничестве с~\href{https://www.overleaf.com/about}{Overleaf} для того, чтобы Вам и Вашим соавторам упростить процесс подготовки рукописей в наш журнал.

Для более подробной информации о журнале обратитесь на его website \url{http://www.mathnet.ru/vsgtu}.

Обратите внимание что в нашем стилевом файле используется команда секционирования \verb"\Section". 
Просим Вас  использовать именно её. 
Если раздел не надо нумеровать, то воспользуетесь вариантом \verb"\Section[N]".

Чтобы подготовить Вашу статью для рецензента, т.~е. обезличить её, воспользуйтесь командой \verb"\ForPeerReview" перед командой 
\verb"\makerutitle".

Помните, что текст статьи  необходимо обязательно структурировать. 
Вы вольны сами выбирать структуру своей статьи.

Обратите внимание, что в нашем журнале необходимо для статей прописывать индексы \href{http://www.mathnet.ru/classifications.phtml?wshow=classifications&option_lang=rus}{УДК и MSC}. 
Просим указывать реально существующие индексы.
Если Вы затрудняетесь в определении индекса УДК, то обратитесь к~любому библиотекарю Вашей научной библиотеки.

\Section{Примеры набора формул и организации ссылок}
Здесь мы приведем некоторые различные наборы формул, которые применялись в~нашем журнале ранее~[\citen{RepKum16}--\citen{{Ung16}}].  
Мы просто из этих статей скопировали небольшие фрагменты в~этот документ.
Рекомендуем использовать окружения {\tt equation}, {\tt multline}, {\tt gather}.


Для набора формул Вы можете пользоваться любым руководством по \LaTeXe.
Мы рекомендуем использовать краткое руководство от К.~В.~Воронцова (\url{http://www.ccas.ru/voron/download/voron05latex.pdf}).

Ссылки на формулы оформляются только командой \verb"\eqref".
Ссылки на таблицы и рисунки оформляются командой \verb"\ref".
Для ссылок на литературу можно использовать команды \verb"\cite" и \verb"\citen".
Обратите внимание, что в нашем журнале нумеруются только те объекты (формулы, рисунки, таблицы, элементы списка литературы), на которые есть ссылки. 
Ссылки на формулы или другие объекты, которые не <<помечены>>, будут обозначаться так ({\bf ??}) или так [{\bf ?}].

В многочисленных исследованиях по наследственной механике, берущих начало с работ Х.~Больцмана \cite{ulg:bib:Boltzman} и продолженных в работах Г.~Дюффинга~\cite{ulg:bib:Duffing}, А.~Джемента \cite{ulg:bib:Gemant1936}, А.~П.~Бронского~\cite{ulg:bib:Bronsky}, 
Г.~Л.~Слонимского~\cite{ulg:bib:Slonimsky}, \linebreak А.~Ю.~Ишлинского~\cite{ulg:bib:Ishlinsky1940}, А.~Р.~Ржаницына~\cite{ulg:bib:Rganitsin}, 
Ю.~Н.~Работнова~\cite{ulg:bib:Rabotnov1948} и других авторов (см. библиографический список в~[\citen{ulg:bib:Bulgakov}--\citen{ulg:bib:Mainardi2010}]), показано и~обосновано, что для целого ряда физических сред ядра ползучести в интегральном операторе \eqref{eq:ulg:VolterraEpsilon} являются ядрами абелевского типа, а основные соотношения между напряжениями и деформациями выражаются через дробные интегралы Римана"--~Лиувилля \cite{ulg:bib:S-K-M}.

Известно, что в рамках структурного моделирования стандартную одномерную обобщённую модель вязкоупругого тела можно записать в виде 
\begin{equation}\label{eq:ulg:GeneralModel}
\sigma(t)+\sum^{n}_{k=1}b_{k}D^{k}\sigma=E\varepsilon+\sum^{m}_{k=1}a_{k}D^{k}\varepsilon,
\end{equation}
где $D^{k}=\left(d/dt\right)^{k}$; $a_{k}$, $b_{k}$, $E_{0}$ "--- постоянные величины, обычно заранее неизвестные; $n=m$ или $m=n+1$. 
Частными случаями модели \eqref{eq:ulg:GeneralModel} являются собственно законы Гука 
$$\sigma(t)=E_{0}\varepsilon(t),$$ и Ньютона $$\sigma(t)=\eta\dot{\varepsilon}(t),$$ а также следующие модели: 
\begin{itemize}
\item модель Фойхта
\begin{equation}\label{eq:ulg:Voigt}
\sigma(t)=E_{0}\varepsilon(t)+\eta\dot{\varepsilon}(t),
\end{equation}
\item модель Максвелла
\begin{equation}\label{eq:ulg:Maxwell}
 E_{0}\sigma(t)+\eta\dot{\sigma}(t)=\eta E_{0}\dot{\varepsilon}(t),
\end{equation}
\item модель Кельвина
\begin{equation}\label{eq:ulg:Kelvin}
\left(E_{1}+E_{2}\right)\sigma(t)+\eta\dot{\sigma}(t)=E_{1}E_{2}\varepsilon(t)+\eta E_{2}\dot{\varepsilon}(t),
\end{equation}
\item модель Зенера
 \begin{equation}\label{eq:ulg:Zener}
 E_{1}\sigma(t)+\eta\dot{\sigma}(t)=E_{1}E_{2}\varepsilon(t)+\eta(E_{1}+E_{2})\dot{\varepsilon}(t),
\end{equation}
\end{itemize}
где $E_{0}$, $E_{1}$, $E_{2}$ "---  модули соответствующих упругих элементов, $\eta$ "---  коэффициент демпфирования. 
Заметим, что всё многообразие реологических моделей можно разбить на два типа "---  модели, описывающие явление мгновенной упругой деформации в момент приложения нагрузки (модели типа Максвелла, Кельвина, Зенера) 
и~модели типа Фойхта, не наделённые (не описывающие) мгновенной упругой деформацией. 
Очевидно, что решения задачи о ползучести при постоянно действующей нагрузке $\sigma(t)=\sigma_{0}=\rm const$, получающиеся из определяющего соотношения \eqref{eq:ulg:GeneralModel}, будут носить в основном экспоненциальный характер. 



Если $\sigma(t)$ изменяется по закону~\eqref{eq:ulg:StressT}, в котором $\sigma_{0}(t)=\sigma_{0}>0$:
\begin{equation}\label{eq:ulg:StressConst}
\sigma(t)=\sigma_{0}\bigl(H(t)-H(t-t_{1})\bigr),
\end{equation}
то
\begin{multline}\label{eq:ulg:ScottBlairSolution}
\varepsilon(t)=\frac{\sigma_{0}}{\eta}I_{0t}^{\alpha} \bigl( H(t)-H(t-t_{1})\bigr)
= 
\\
=\frac{\sigma_{0}}{\eta\Gamma(\alpha)} \biggl(H(t)\int _{0}^{t}(t-\tau)^{\alpha-1}d\tau
-H(t-t_{1}) \biggl(\int _{0}^{t_{1}}+\int _{t_{1}}^{t} \biggr)(t-\tau)^{\alpha-1}d\tau \biggr)=
\\
= \frac{\sigma_{0}}{\eta \Gamma (\alpha+1)} \bigl(t^{\alpha}H(t)-(t-t_{1})^{\alpha}H(t-t_{1})\bigr).
\end{multline}

Отметим, что при $\alpha\rightarrow 0$ решение~\eqref{eq:ulg:ScottBlairSolution} сводится к
закону Гука, если положить $\eta=E_{0}$, а при $\alpha\rightarrow 1$ определяет линейный
закон деформирования.

\Section{Пользуйтесь окружениями для теорем, лемм, задач и прочего}
Вы можете использовать следующие определенные в пакете {\tt samgtu} окружения {\tt definition}, {\tt lemma}, {\tt example}, {\tt remark}, {\tt task}, {\tt proof}.

\begin{lemma}
Это текст леммы.
\end{lemma}

\begin{remark}[1]
Это текст замечания. Номер окружениям  {\tt definition}, {\tt lemma}, {\tt example}, {\tt remark}, {\tt task} проставляется пока вручную.
\end{remark}



Рассмотрим уравнение
\begin{equation}
       \label{repin:eq1}
0= \left\{\begin {array}{ll}
      u_{xxx}-u_{y}, & y>0,\\
      y^{2m}u_{xx}+y\,u_{yy}+\alpha u_{y}, & y<0,
\end{array} \right.
\end{equation}
где $m$ "--- натуральное число, 
$\alpha=\text{const}$, $({1-2m})/{2}<\alpha<1 $ в конечной области $\Omega$, 
ограниченной отрезками $AA_{0}$, $BB_{0}$, $A_{0}B_{0}$  прямых $x=0$, $x=1$, $y=1$ соответственно и характеристиками
\[
AC: \, x-\frac{2}{2m+1}(-y)^{\frac{2m+1}{2}}=0, \quad BC: \, x+\frac{2}{2m+1}(-y)^{\frac{2m+1}{2}}=1
\]
уравнения \eqref{repin:eq1} при $y<0$.


\begin{newthm}{Теорема} В области $\Omega$ не может существовать более одного решения задачи \eqref{repin:eq1}--\eqref{repin:eq4} при $({1-2m})/{2}< \alpha<1,$ если либо
\begin{gather}
       \label{repin:eq5}
\alpha_{1}=\beta_{1}=1-\beta, \quad  w(x)=\delta(x) \equiv 1,
\\
       \label{repin:eq8}
M_{1}(x)=\gamma_{1}(1-x)^{\beta} a(x)+\gamma_{1}x^{\beta} b(x)-\frac{\Gamma(\beta)}{\Gamma(2\beta)} x^{\beta} (1-x)^{\beta}d(x) \ne 0,
\\
       \label{repin:eq8}
\left[ \frac{(1-x)^{\beta} a(x)}{M_{1}(x)} \right]' \leq 0, 
\;
\left[ \frac{x^{\beta} b(x)}{M_{1}(x)} \right]' \geq 0, 
\; 
\frac{x^{\beta}(1-x)^{\beta} c(x)}{M_{1}(x)} \geq 0, 
\;
\forall x \in \overline I,
 \end{gather}
либо
\begin{equation}
       \label{repin:eq8}
\alpha_{1}=\beta_{1}=\beta, \quad \delta(x) =x^{2\beta-1}, \quad w(x)=(1-x)^{2\beta-1}
\end{equation}
и выполняются условия
\begin{gather}
       \label{repin:eq10}
M_{2}(x)=(1-x)^{1-\beta} a(x)+x^{1-\beta} b(x)+\frac{\Gamma(\beta)}{\Gamma(2\beta)} x^{1-\beta} (1-x)^{1-\beta}c(x) \ne 0,
\\
       \label{repin:eq10}
       \!\!\!\!\!\! \hspace{-3mm}
\left[ \frac{\gamma_{1}(1-x)^{\beta} a(x)}{M_{2}(x)} \right]’ \leq 0,  
\left[ \frac{\gamma_{1} x^{1-\beta} b(x)}{M_{2}(x)} \right]’ \geq 0, 
\frac{x^{1-\beta}(1-x)^{1-\beta} d(x)}{M_{2}(x)} \leq 0,  \forall x \in \overline I,
 \end{gather}
где
\[
\gamma_{1}=\frac{\Gamma(\beta) \Gamma(1-2\beta)}{2\Gamma(2 \beta) \Gamma(1-\beta)} \Bigl( \frac{2m+1}{4}\Bigr)^{-2\beta}.
\]
\end{newthm}

\smallskip

\begin{proof}
Регулярное в области $\Omega_{2}$ решение уравнения \eqref{repin:eq1} при $({1-2m})/{2}< \alpha<1$, удовлетворяющее условиям
\begin{equation}
       \label{repin:eq11}
u(x, 0)= \tau(x), \quad \lim_{y \to -0} (-y)^{\alpha}u_{y}= \nu(x), \quad 0<x<1,
\end{equation}
единственно и имеет вид [{\citen{repin:7}]
\begin{multline}
       \label{repin:eq12}
u(x,y)= \frac{\Gamma(2\beta)}{\Gamma^{2}(\beta)}  \int _{0}^{1} \tau \left[ x+ \frac{2(1-2t)}{2m+1} (-y)^{\frac{2m+1}{2}}\right]t^{\beta-1}(1-t)^{\beta-1}dt- \\
- \frac{2}{m+1} \frac{\Gamma(1-2\beta)}{\Gamma^{2}(1-\beta)} (-y)^{1-\alpha}  \int _{0}^{1} \nu \left[ x+ \frac{2(1-2t)}{2m+1} (-y)^{\frac{2m+1}{2}}\right]t^{-\beta}(1-t)^{-\beta}dt,
\end{multline}
где
\[
\beta=\frac{2m-1+2\alpha}{2(2m+1)}.
\]

При выполнении условий \eqref{repin:eq5} теоремы соотношение между $\tau(x)$ и $\nu(x)$ из области $\Omega_{2}$ запишем в~виде
\begin{equation}
       \label{repin:eq13}
\nu(x) =A_{1}(x) D_{0x}^{1-2\beta} \tau(x)+ B_{1}(x) D_{x1}^{1-2\beta} \tau(x)+ C_{1}(x) \tau(x) +F_{1}(x),
\end{equation}
где
\[
A_{1}(x)=\frac{(1-x)^{\beta} a(x)}{M_{1}(x)}, \quad B_{1}(x)=\frac{x^{\beta} b(x)}{M_{1}(x)},
\]
\[
C_{1}(x)=\frac{x^{\beta}(1-x)^{\beta} c(x)}{M_{1}(x)}, \quad F_{1}(x)=-\frac{\Gamma(\beta)}{\Gamma(2\beta)} \, \frac{x^{\beta}(1-x)^{\beta} \gamma(x)}{M_{1}(x)}.
\]



Аналогичными рассуждениями из \eqref{repin:eq15} можно получить $\tau_{i}(x)=0$, $i=1, 2$ и, следовательно, $\tau(x)=0$, а из \eqref{repin:eq13} при $\gamma(x)=0$ имеем $\nu(x)=0$. 
Таким образом, решение задачи $u(x,y) \equiv $ в $\Omega_{2}$ как решение задачи \eqref{repin:eq11} с нулевыми данными, а в области $\Omega_{1}$ как решение однородной задачи \eqref{repin:eq1}, \eqref{repin:eq2}.
}
\end{proof}

\Section{Вставляйте рисунки и таблицы} 
Обратите внимание, что  данные на рисунках и в таблицах необходимо переводить на английский язык. 

Вставляются рисунки и таблицы с помощью окружений {\tt figure} и {\tt table}. 
Обратите внимание, что рисунки размещаются в каталоге {\tt ./00PIC/}.
Рисунки должны быть обязательно векторными. 
Только такие рисунки мы можем подправить  и внедрить в них необходимые шрифты.
Возможно использовать растровые рисунки для фотографий, скриншотов, но такие рисунки должны быть оригинальными, они не должны быть пережаты.

В результате первичной обработки каждой кривой вязкоупругого деформирования при всех пяти напряжениях $\sigma_{0}$
при нагрузке
 $$\sigma(t)=\sigma_{0}={\rm const}~~ (0\leq t \leq 8),\quad \sigma(t)=0~~ (8<t\leq 12)$$ 
 найдены значения параметров для всех рассмотренных моделей.
 
\begin{table}[b!]

\small \centering

\caption{Значения параметров аппроксимации~(\ref{FracKelvinSolution}) и погрешности для дробного аналога модели Кельвина [The values of the approximation parameters for Eq.~(\ref{FracKelvinSolution}) and measure of inaccuracy for fractional Kelvin model] \label{un:tabl1}}
\vspace{-3mm}
\renewcommand{\tabcolsep}{0.55cm}
\begin{tabular}{c||c|c|c|c|c}
\hline
\rule{0mm}{11pt}%
\footnotesize $\sigma_{0}$, MPa & \footnotesize  $\alpha$ & \footnotesize  $E_{1}$, MPa & \footnotesize  $E_{2}$, MPa & \footnotesize  $\eta$ & \footnotesize  $\Delta,\, \%$ \\
            \hline
            \rule{0mm}{12pt}%
            \phantom{1}4.655 & 0.355 & 145.525 & 1163.800 & 190.067 & 2.586 \\
             \phantom{1}6.288 & 0.326 & 133.601 & \phantom{1}898.329 & 158.793 & 4.097 \\
             \phantom{1}8.738 & 0.394 & 140.194 & 1028.000 & 128.956 & 3.168 \\
             10.372 & 0.318 & \phantom{1}97.571 & \phantom{1}829.733 & 158.840 & 3.468 \\
             12.005 & 0.400 & 127.962 & \phantom{1}857.500 & 113.856 & 3.121 \\
            \hline
        \end{tabular}

\end{table} 
 

В табл.~\ref{un:tabl1} приведены значения параметров аппроксимации для дробного аналога модели Кельвина. 
В~колонке $\Delta$ с использованием нормы~\eqref{norma}  приведены отклонения теоретической зависимости $\varepsilon(t_{i})$,  рассчитанной по формуле~\eqref{FracKelvinSolution} при каждом значении напряжения $\sigma_{0}$ с найденными параметрами $\alpha$, $E_{1}$, $E_{2}$, $\eta$, от экспериментальных данных. 
В~качестве примера на рис.~\ref{Kelvin570g} штриховой линией представлена полученная по формуле~\eqref{FracKelvinSolution} расчетная зависимость при $\sigma_{0}=4.655$~MПa с~параметрами из первой строки табл.~\ref{un:tabl1}. 
Погрешность аппроксимации конкретной реализации в данном случае равна 2.586\,\%. 
Анализ данных табл.~\ref{un:tabl1} свидетельствует о том, что, вообще говоря, для дробного аналога модели Кельвина наблюдается существенный разброс параметров модели для реализаций при различных значениях напряжений.

\begin{figure}[h!]
\noindent
\begin{minipage}[h]{0.48\textwidth}

\caption{Экспериментальная (сплошная линия) и расчетная по модели~(\ref{FracKelvinSolution}) (штриховая линия) кривые вязкоупругого деформирования поливинилхлоридного пластиката при напряжении $\sigma_{0}=4.655$~МПа с~последующей разгрузкой
\label{Kelvin570g}
 }

\smallskip \footnotesize

[Figure~\ref{Kelvin570g}.
Experimental (solid line) and calculated by the  Kelvin's fractional  mo\-del~\eqref{FracKelvinSolution} (dashed line) viscoelastic creep curves of the flexible PVC at the stress $\sigma_{0}=4.655$~MPa with subsequent \centerline{unloading]}


\end{minipage} \hfill
\begin{minipage}[h]{0.5\textwidth}


\centering
\includegraphics[width=0.98\textwidth]{Kelvin570g}

\end{minipage}

\end{figure}


Аналогичный расчет проведен и для других рассмотренных в настоящей работе моделей при всех пяти уровнях напряжений. Для построения исследуемых моделей, <<работоспособных>> при напряжениях $\sigma_{0}$ от $4.655$~МПа до $ 12.005$~МПа, выполнено усреднение параметров по пяти реализациям для каждой из них.  В табл.~\ref{un:tabl2} представлены полученные значения усредненных параметров $\bar{\alpha}$, $\bar{E}_{1}$, $\bar{E}_{2}$, $\bar{\eta}$ для всех рассмотренных моделей.




\begin{figure}[p!]

\noindent
\begin{minipage}[h]{0.5\textwidth}
\centering  \scriptsize \sl
\includegraphics[width=0.98\textwidth]{Kelvin570g-a}\\
a
\end{minipage}
\begin{minipage}[h]{0.5\textwidth}
\centering  \scriptsize \sl
\includegraphics[width=0.98\textwidth]{Kelvin770g-b}\\
b
\end{minipage}

\vspace{2mm}

\noindent
\begin{minipage}[h]{0.5\textwidth}
\centering  \scriptsize \sl
\includegraphics[width=0.98\textwidth]{Kelvin1070g-c}\\
c
\end{minipage}
\begin{minipage}[h]{0.5\textwidth}
\centering  \scriptsize \sl
\includegraphics[width=0.98\textwidth]{Kelvin1270g-d}\\
d
\end{minipage}

\vspace{2mm}

\begin{minipage}[h]{0.5\textwidth}
\centering  \scriptsize \sl
\includegraphics[width=0.98\textwidth]{Kelvin1470g-e}\\
e
\end{minipage} \hfill
\begin{minipage}[h]{0.48\textwidth}
\caption{Экспериментальные (сплошные линии) и расчетные по модели~(\ref{FracKelvinSolution}) с усредненными значениями параметров
(штриховые линии) кривые вязкоупругого деформирования поливинилхлоридного пластиката при напряжениях $\sigma_{0}=4.655$~МПа ({\sl a}\/),
$\sigma_{0}=6.288$~МПа ({\sl b}\/),
$\sigma_{0}=8.738$~МПа ({\sl c}\/), $\sigma_{0}=10.372$~МПа ({\sl d}\/) и
$\sigma_{0}=12.005$~МПа ({\sl e}\/)
 с~последующей разгрузкой \label{KelvinAll} }
\end{minipage}

\smallskip \footnotesize

[Figure~\ref{KelvinAll}.
Experimental (solid line) and calculated by the Kelvin's fractional  mo\-del (\ref{FracKelvinSolution})   with averaged values of parameters  (dashed line) viscoelastic creep curves of the flexible PVC at the stresses \mbox{$\sigma_{0}=4.655$~MPa} ({\sl a}\/), ${\sigma_{0}=6.288}$~MPa ({\sl b}\/),
$\sigma_{0}=8.738$~MPa ({\sl c}\/), $\sigma_{0}=10.372$~MPa ({\sl d}\/), \centerline{ and
$\sigma_{0}=12.005$~MPa ({\sl e}\/) with subsequent unloading]}
\end{figure}


\begin{table}[t!]

\small \centering

\caption{Усредненные значения параметров для дробных моделей \label{un:tabl2}} \vspace{-3mm}
[The averaged parameter values for the fractional models] 
\renewcommand{\tabcolsep}{0.3cm}
\begin{tabular}{l||c|c|c|c} %\label{tab:ModelsParameters}
\hline
\multicolumn{1}{c||}{\footnotesize Fractional Models} & \footnotesize  $\bar{\alpha}$ & \footnotesize  $\bar{E}_{1}$, MPa & \footnotesize  $\bar{E}_{2}$, MPa & \footnotesize  $\bar{\eta}$\rule{0mm}{11pt}\\
            \hline
            \rule{0mm}{12pt}%
           The Scott Blair's fractional model (\ref{SB_Solution}) & 0.179  &  ---  &  ---  & 289.847   \\
           The Voigt's fractional model  (\ref{FracVoightSolution}) & 0.359  & \phantom{1}128.971 & --- &  150.102 \\
           The Maxvell's fractional model  (\ref{FracMaxwellSolution})  & 0.179 &  ---  & 955.472 & 289.847 \\
           The Kelvin's fractional  model (\ref{FracKelvinSolution}) & 0.359 & \phantom{1}128.971 & 955.472 & 150.102 \\
           The Zener's fractional  model  (\ref{FracZenerSolution}) & 0.285 & 1078.692 & \phantom{1}92.418 & 143.029 \\
            \hline
        \end{tabular}

\bigskip

\caption{Погрешности аппроксимаций ($\Delta, \%$) для всех исследуемых моделей после осреднения параметров для пяти реализаций
[The approximation errors $\Delta$ (in percentages) for all investigated models after averaging parameters for the five realizations] \label{un:tabl3}} 

\vspace{-3mm}
\renewcommand{\tabcolsep}{0.5cm}
\begin{tabular}{l||c|c|c|c|l} %\label{tab:ModelsError}
\hline
\rule{0mm}{11pt}%
 & \multicolumn{5}{c}{Stresses $\sigma_0$, MPa}\\ \cline{2-6}
\totop{\footnotesize Fractional Models} & \footnotesize  4.655  & \footnotesize  6.288 & \footnotesize  8.738  & \footnotesize  10.372 & \footnotesize  12.005  \\
\hline
\multicolumn{1}{l||}{The model (\ref{SB_Solution})} & 19.139  &  12.538  &  10.192  &  13.237  &  15.745\rule{0mm}{12pt}\\
\multicolumn{1}{l||}{The model  (\ref{FracVoightSolution})} & 18.580  &  \phantom{1}8.892  &  \phantom{1}4.841  &  11.062  &  12.273 \\
\multicolumn{1}{l||}{The model  (\ref{FracMaxwellSolution})} & 19.792  &  \phantom{1}9.671  &  \phantom{1}8.578  &  12.609  &  14.429 \\
\multicolumn{1}{l||}{The model (\ref{FracKelvinSolution})} & 19.325  &  \phantom{1}6.606  &  \phantom{1}3.964  &  11.153  &  11.843 \\
\multicolumn{1}{l||}{The model  (\ref{FracZenerSolution})} & 19.018  &  \phantom{1}7.394  &  \phantom{1}4.937  &  11.438  &  12.527 \\
\hline
        \end{tabular}


\end{table}





В табл.~\ref{un:tabl3}  приведены погрешности отклонений расчетных значений вязкоупругой деформации от экспериментальных для всех пяти моделей с усредненными параметрами, вычисленные по формуле~\eqref{norma} при каждом уровне напряжений. Средняя же погрешность по всем пяти реализациям для модели Скотт Блэра равна 14.170\,\%, для дробных аналогов модели Фойхта     "---     11.130\,\%, Максвелла     "---     13.016\,\%, Кельвина     "---     10.578\,\% и Зенера     "---     11.063\,\%.
Таким образом, наилучшее приближение теоретической кривой к экспериментальным значениям дает дробный аналог модели Кельвина. В качестве иллюстрации на рис.~\ref{KelvinAll} и~\ref{MaxwellAll} штриховой линией построены расчетные значения вязкоупругой деформации по моделям~\eqref{FracKelvinSolution} и~\eqref{FracMaxwellSolution} соответственно с усредненными коэффициентами.

\Section{Немного про библиографический список}
Для оформления библиографического списка в нашем журнале используется пакет {\tt samgtu-bib}. 
Этот пакет "--- модификация пакета {\tt amsbib}, который размещён на сайте \href{http://www.mathnet.ru/poffice/amsbibpackage.phtml?wshow=amsbibpackage&option_lang=rus}{Math-Net.Ru} и  используется в~ведущих журналах Отделения математических наук РАН. 
Поэтому для пакета {\tt samgtu-bib} верны все примеры, которые приводятся в руководстве к пакету {\tt amsbib} (\url{http://www.mathnet.ru/poffice/amsbib.pdf}).

Просим Вас приводить в библиографическом списке выверенную информацию.
Не приводите источники, которые трудно найти, например, тезисы конференций, препринты, не имеющие электронного варианта в сети интернет.
Также не рекомендуем приводить ссылки на web-ресурсы, не являющиеся научными журналами.

Формируйте один библиографический список.
Члены редколлегии разобьют его на два самостоятельно.


Старайтесь в библиографическом списке приводить оригинальную, не переводную литературу.
Если русскоязычные источники имеют  перевод на английский язык, или у них есть официальные meta-данные на английском языке, то приводите их.
Посмотрите как это делается в библиографическом списке здесь.

Если Вы ссылаетесь на статью, размещенную на портале \href{http://www.mathnet.ru/}{Math-Net.Ru}, то смело может копировать её цитирование в формате  {\tt amsbib} с~портала.
Если Вы знаете \DOIlogo\ для статей из вашего библиографического списка, то обязательно приводите их.


\Section[N]{Заключение (Выводы)}
Каждая статья должна заканчиваться заключением или выводами по работе. 
Приведите здесь краткую формулировку результатов своего исследования. 
После этого раздела (перед библиографическим списком) Вы можете выразить благодарности своим коллегам или предоставить другую информацию с помощью команды \verb"\AdditionalContent".
Обратите внимание, что дополнительная информация обязательно переводится на английский язык.
Примеры такой дополнительной информации приведены ниже.
Посмотрите какую дополнительную информацию размещают наши авторы, авторы в других журналах.

Не забудьте сгенерировать страницу с информацией о статье на английском языке с помощью команды \verb"\makeentitle".




\AdditionalContent{Конкурирующие интересы}{Укажите какие конфликты интересов имеются.} 
%\AdditionalContent{Конкурирующие интересы}{Конкурирующих интересов не имею.} 
\AdditionalContent{Авторский вклад и ответственность}{Укажите степень авторской ответственности каждого автора.} 
%\AdditionalContent{Авторская ответственность}{Я несу полную ответственность за предоставление окончательной версии рукописи в печать. Окончательная версия рукописи мною одобрена.} 
%\AdditionalContent{Авторский вклад и ответственность}{Все  авторы  принимали  участие  в  разработке  концепции  \mbox{статьи}  и  в  написании  рукописи. Авторы несут  полную  ответственность  за  предоставление  окончательной рукописи в печать. Окончательная  версия рукописи была одобрена всеми авторами.} 
\AdditionalContent{Финансирование}{Укажите источник финансирования работы.}
%\AdditionalContent{Финансирование}{Исследование выполнялось без финансирования.}
%\AdditionalContent{Финансирование}{Работа выполнена при поддержке РФФИ (проект № xx–xx–xxxxx_a).}
\AdditionalContent{Благодарность}{Выразите свою благодарность.}
%\AdditionalContent{Благодарность}{Авторы благодарны рецензенту за тщательное прочтение статьи и~ценные предложения и комментарии.}

\begin{thebibliography}{10}

\RBibitem{RepKum16}
\by Репин~О.~А., Кумыкова~С.~К.
\paper Внутреннекраевая задача с операторами Римана--Лиувилля для уравнения смешанного типа~третьего порядка
\jour Вестн. Сам. гос. техн. ун-та. Сер. Физ.-мат. науки
\yr 2016
\vol 20
\issue 1
\pages 43--53
\crossref{10.14498/vsgtu1461}
\misctransl
\lang In Russian
\by Repin~O.~A., Kumykova~S.~K.
\paper An internal boundary value problem with~the Riemann--Liouville operator for the mixed type equation of the third order
\jour Vestn. Samar. Gos. Tekhn. Univ. Ser. Fiz.-Mat. Nauki \rm [\ParallelJournalName]
\yr 2016
\vol 20
\issue 1
\pages 43--53

\RBibitem{OgoRadUng16}
\by Огородников~Е.~Н., Радченко~В.~П., Унгарова~Л.~Г.
\paper Математическое моделирование наследственно упругого деформируемого тела 
на основе структурных моделей и~аппарата дробного 
интегро-дифференцирования Римана--Лиувилля
\jour Вестн. Сам. гос. техн. ун-та. Сер. Физ.-мат. науки
\yr 2016
\vol 20
\issue 1
\pages 167--194
\crossref{10.14498/vsgtu1456}
\misctransl
\lang In Russian
\by Ogorodnikov~E.~N., Radchenko~V.~P., Ungarova~L.~G.
\paper Mathematical modeling of hereditary elastically deformable body on the basis 
of structural models and fractional integro-differentiation Riemann--Liouville apparatus
\jour Vestn. Samar. Gos. Tekhn. Univ. Ser. Fiz.-Mat. Nauki \rm [\ParallelJournalName]
\yr 2016
\vol 20
\issue 1
\pages 167--194


\RBibitem{Ung16}
\by Унгарова~Л.~Г.
\paper Применение линейных дробных аналогов реологических моделей в~задаче аппроксимации экспериментальных данных 
по растяжению поливинилхлоридного пластиката
\jour Вестн. Сам. гос. техн. ун-та. Сер. Физ.-мат. науки 
\yr 2016
\vol 20
\issue 4
\pages 691--706
\crossref{10.14498/vsgtu1523}
\misctransl
\lang In Russian
\by Ungarova~L.~G.
\paper The use of linear fractional analogues of rheological models in the problem of~approximating the experimental data 
on~the stretch polyvinylchloride elastron
\jour Vestn. Samar. Gos. Tekhn. Univ. Ser. Fiz.-Mat. Nauki \rm [\ParallelJournalName]
\yr 2016
\vol 20
\issue 4
\pages 691--706


\RBibitem{ulg:bib:Rabotnov1966}
\by Работнов~Ю.~Н.
\book Ползучесть элементов конструкций
\yr 1966
\publ Наука
\publaddr М.
\totalpages 752
\transl
\by Rabotnov~Yu.~N.
\book Creep problems in structural members
\yr 1969
\publ North-Holland Publ. Co.
\publaddr Amsterdam, London
\totalpages xiv+822

\zmath{https://zbmath.org/?q=an:03294380}

\RBibitem{ulg:bib:RabotnovMDTT}
\by Работнов~Ю.~Н.
\book Механика деформируемого тв\"eрдого тела
\yr 1988
\publ Наука
\publaddr М.
\totalpages 712
\zmath{https://zbmath.org/?q=an:00193379}
\misctransl
\by Rabotnov~Yu.~N.
\book Mekhanika deformiruemogo tverdogo tela \rm [Mechanics of a deformable rigid body]
\yr 1988
\publ Nauka
\publaddr Moscow
\lang In Russian
\totalpages 712


\RBibitem{ulg:bib:malinin}
\by Малинин~Н.~Н.
\book Прикладная теория пластичности и ползучести
\yr 1975
\publ Машиностроение
\publaddr М.
\totalpages 400
\misctransl
\by Malinin~N.~N.
\book Prikladnaia teoriia plastichnosti i polzuchesti \rm [Applied theory of plasticity and creep]
\yr 1975
\publ Mashinostroenie
\publaddr Moscow
\lang In Russian
\totalpages 400


\RBibitem{ulg:bib:Rabotnov1977}
\by Работнов~Ю.~Н.
\book Элементы наследственной механики тв\"eрдых тел
\yr 1977
\publ Наука.
\publaddr М.
\totalpages 384
\transl
\book Elements of hereditary solid mechanics
\yr 1980
\publ Mir Publ.
\totalpages 388
\zmath{https://zbmath.org/?q=an:0515.73026}

\RBibitem{ulg:bib:NahushevFracCalculus}
\by Нахушев~А.~М.
\book Дробное исчисление и его применение
\yr 2003
\publ Физматлит
\publaddr М.
\totalpages 272
\zmath{https://zbmath.org/?q=an:1066.26005}
\misctransl
\by Nakhushev~A.~M.
\book Drobnoe ischislenie i ego primenenie \rm [Fractional calculus and its applications]
\yr 2003
\publ Fizmatlit
\publaddr Moscow
\lang In Russian
\totalpages 272


\RBibitem{ulg:bib:NahushevMatBiology}
\by Нахушев~А.~М.
\book Уравнения математической биологии
\yr 1995
\publ Высшая школа
\publaddr М.
\totalpages 301
\zmath{https://zbmath.org/?q=an:0991.35500}
\misctransl
\by Nakhushev~A.~M.
\book Uravneniia matematicheskoi biologii \rm [Equations of mathematical biology]
\yr 1995
\publ Vysshaia shkola
\publaddr Moscow
\lang In Russian
\totalpages 301


\Bibitem{ulg:bib:Uchaikin2013}
\by Uchaikin~V.~V.
\paper Heredity and Nonlocality
\inbook Fractional Derivatives for Physicists and Engineers
\yr 2013
\pages 3--58
\publ Springer Berlin Heidelberg
\publaddr Berlin
\bookvol 1
\volinfo Background and Theory
\crossref{10.1007/978-3-642-33911-0_1}
\adsnasa{http://adsabs.harvard.edu/cgi-bin/bib_query?2013fdpe.book.....U}
\moreref
\crossref{10.1007/978-3-642-33911-0}

\Bibitem{ulg:bib:Volterra1909}
\by Volterra~V.
\paper Sulle equazioni integro-differenziali della teoria dell'elasticit\`a
\jour Rend. Acc. Naz. Lincei
\yr 1909
\vol 18
\pages 295–301


\RBibitem{ulg:bib:Volterra1982}
\by Вольтерра~В.
\book Теория функционалов, интегральных и интегро-дифференциальных уравнений
\yr 1982
\publ Наука
\publaddr М.
\totalpages 304
\transl
\by Volterra~V.
\book Theory of functionals and of integral and integro-differential equations
\yr 1959
\publ Dover Publ., Inc.
\publaddr New York
\totalpages 226
\zmath{https://zbmath.org/?q=an:0086.10402|02573174}

\Bibitem{ulg:bib:Boltzman}
\by Boltzmann~L.
\paper Theorie der elastischen Nachwirkung (Theory of elastic after effects)
\jour Wien. Ber.
\yr 1874
\vol 70
\pages 275–306
\zmath{https://zbmath.org/?q=an:02717441}
\moreref
\by Boltzman~L.
\paper Zur Theorie der elastischen Nachwirkung (On the elastic after effect)
\jour Pogg. Ann. \rm (2)
\yr 1878
\vol 5
\pages 430--432
\zmath{https://zbmath.org/?q=an:10.0670.01}
\moreref
\by Boltzman~L.
\paper Zur Theorie der elastischen Nachwirkung
\inbook Wissenschaftliche Abhandlungen
\serial Cambridge Library Collection
\yr 2012
\pages 318--320
\publ Cambridge University Press
\publaddr Cambridge
\ed Friedrich Hasen\"ohrl
\bookvol 2
\crossref{10.1017/CBO9781139381437.015}

\Bibitem{ulg:bib:Duffing}
\by Duffing~G.
\paper Elastizit\"at und Reibung beim Riementrieb (Elasticity and friction of the belt drive)
\jour Forschung auf dem Gebiet des Ingenieurwesens~A
\yr 1931
\vol 2
\issue 3
\pages 99--104
\crossref{10.1007/BF02578795}
\scopus{http://www.scopus.com/record/display.url?origin=inward&eid=2-s2.0-51249193154}

\Bibitem{ulg:bib:Gemant1936}
\by Gemant~A.
\paper A Method of Analyzing Experimental Results Obtained from Elasto-Viscous Bodies
\jour J. Appl. Phys.
\yr 1936
\vol 7
\pages 311--317
\crossref{10.1063/1.1745400}
\adsnasa{http://adsabs.harvard.edu/cgi-bin/bib_query?1936JAP.....7..311G}
\scopus{http://www.scopus.com/record/display.url?origin=inward&eid=2-s2.0-0345218900}

\RBibitem{ulg:bib:Bronsky}
\by Бронский~А.~П.
\paper Явление последействия в тв\"eрдом теле
\jour ПММ
\yr 1941
\vol 5
\issue 1
\pages 31--56
\zmath{https://zbmath.org/?q=an:0063.00617}
\misctransl
\by Bronskij~A.~P.
\paper Residual effect in rigid bodies
\jour Prikl. Mat. Mekh.
\yr 1941
\vol 5
\issue 1
\pages 31--56
\lang In Russian

\RBibitem{ulg:bib:Slonimsky}
\by Слонимский~Г.~Л.
\paper О законах деформации реальных материалов
\jour ЖТФ
\yr 1939
\vol 9
\issue 20
\pages 1791--1799
\transl
\by Slonimsky~G.~L.
\paper On the laws of deformation of real materials. I. The theories of Maxwell and Boltzmann
\jour Acta physicochim. URSS
\yr 1940
\vol 12
\pages 99--128
\zmath{https://zbmath.org/?q=an:66.1023.04}

\RBibitem{ulg:bib:Ishlinsky1940}
\by Ишлинский~А.~Ю.
\paper Об уравнениях пространственного деформирования не вполне упругих и~вязкопластических тел
\jour Изв. АН СССР, ОТН
\yr 1945
\issue 3
\pages 250–260
\zmath{https://zbmath.org/?q=an:0063.02988}
\misctransl
\by Ishlinsky~A.~Yu.
\paper On equations of spatial deformation of not completely elastic and elastoplastic bodies
\jour Izv. AN SSSR, OTN
\yr 1945
\issue 3
\pages 250–260
\lang In Russian

\RBibitem{ulg:bib:Rganitsin}
\by Ржаницын~А.~Р.
\book Некоторые вопросы механики систем, деформирующихся во времени
\yr 1949
\publ Гостехиздат
\publaddr М.
\totalpages 248
\misctransl
\by Rzhanitsyn~A.~R.
\book Nekotorye voprosy mekhaniki sistem, deformiruiushchikhsia vo vremeni \rm [Some Problems in the Mechanics of Time-Deformable Systems]
\yr 1949
\publ Gostekhizdat
\publaddr Moscow
\lang In Russian
\totalpages 248


\RBibitem{ulg:bib:Rabotnov1948}
\by Работнов~Ю.~Н.
\paper Равновесие упругой среды с последействием
\jour ПММ
\yr 1948
\vol 12
\issue 1
\pages 53--62
\zmath{https://zbmath.org/?q=an:03055326}
\transl
\by Rabotnov~Yu.~N.
\paper Equilibrium of an elastic medium with after-effect
\jour Fractional Calculus and Applied Analysis
\yr 2014
\vol 17
\issue 3
\pages 684--696
\crossref{10.2478/s13540-014-0193-1}
\zmath{https://zbmath.org/?q=an:1306.74011}
\scopus{http://www.scopus.com/record/display.url?origin=inward&eid=2-s2.0-84904339295}

\RBibitem{ulg:bib:Bulgakov}
\by Булгаков~И.~И.
\book Ползучесть полимерных материалов
\yr 1973
\publ Наука
\publaddr М.
\totalpages 288
\misctransl
\lang In Russian
\by Bulgakov~I.~I.
\book Polzuchest' polimernykh materialov \rm [Creep of Polymer Materials]
\yr 1973
\publ Nauka
\publaddr Moscow
\totalpages 288


\Bibitem{ulg:bib:Podlubny1999}
\by Podlubny~I.
\book Fractional differential equations. An introduction to fractional derivatives, fractional differential equations, to methods of their solution and some of their applications
\serial Mathematics in Science and Engineering
\yr 1999
\vol 198
\publ Academic Press
\publaddr San Diego
\totalpages xxiv+340
\zmath{https://zbmath.org/?q=an:01329842}

\Bibitem{ulg:bib:K-S-T}
\by Kilbas~A.~A., Srivastava~H.~M., Trujillo~J.~J.
\book Theory and Applications of Fractional Differential Equations
\serial North-Holland Mathematics Studies
\yr 2006
\vol 204
\publ Elsevier
\publaddr Amsterdam
\totalpages xx+523
\zmath{https://zbmath.org/?q=an:02217537}

\Bibitem{ulg:bib:Mainardi2010}
\by Mainardi~F.
\book Fractional calculus and waves in linear viscoelasticity. An introduction to mathematical models
\yr 2010
\publ World Scientific
\publaddr Hackensack, NJ
\totalpages xx+347
\crossref{10.1142/9781848163300}
\zmath{https://zbmath.org/?q=an:1210.26004}

\RBibitem{ulg:bib:S-K-M}
\by Самко~С.~Г. Килбас~А.~А., Маричев~О.~И.
\book Интегралы и производные дробного порядка и некоторые их приложения
\yr 1987
\publ Наука и техника
\publaddr Минск
\totalpages 688
\zmath{https://zbmath.org/?q=an:0617.26004}
\misctransl
\by Samko~S.~G. Kilbas~A.~A., Marichev~O.~I.
\book Integraly i proizvodnye drobnogo poriadka i nekotorye ikh prilozheniia \rm [Integrals and derivatives of fractional order and some of their applications]
\yr 1987
\publ Nauka i tekhnika
\publaddr Minsk
\lang In Russian
\totalpages 688


\Bibitem{ulg:bib:Gemant1938}
\by Gemant~A.
\paper On fractional differentials
\jour Philos. Mag., VII. Ser.
\yr 1938
\vol 25
\pages 540--549
\crossref{10.1080/14786443808562036}
\zmath{https://zbmath.org/?q=an:0018.25201}

\RBibitem{ulg:bib:Ogorodnikov_Some_Special_Functions}
\by Огородников~Е.~Н., Яшагин~Н.~С.
\paper Некоторые специальные функции в решении задачи Коши для одного дробного осцилляционного уравнения
\jour Вестн. Сам. гос. техн. ун-та. Сер. Физ.-мат. науки
\yr 2009
\issue 1(18)
\pages 276--279
\mathnet{http://mi.mathnet.ru/rus/vsgtu685}
\crossref{10.14498/vsgtu685}
\misctransl
\by Ogorodnikov~E.~N., Yashagin~N.~S.
\paper Some Special Functions in the Solution To Cauchy Problem for a Fractional Oscillating Equation
\jour Vestn. Samar. Gos. Tekhn. Univ. Ser. Fiz.-Mat. Nauki \rm [J. Samara State Tech. Univ., Ser. Phys. \& Math. Sci.]
\yr 2009
\issue 1(18)
\pages 276--279
\lang In Russian


\RBibitem{ulg:bib:Djrbashyan1966}
\by Джрбашян~М.~М.
\book Интегральные преобразования и представления функций в~комплексной области
\yr 1966
\publ Наука
\publaddr М.
\totalpages 672
\zmath{https://zbmath.org/?q=an:0154.37702}
\misctransl
\by Dzhrbashyan~M.~M.
\book Integral'nye preobrazovaniia i predstavleniia funktsii v~kompleksnoi oblasti \rm [Integral Transforms and Representation of Functions in Complex Domain]
\yr 1966
\publ Nauka
\publaddr Moscow
\lang In Russian
\totalpages 672


\RBibitem{ulg:bib:ORY_reology_models}
\by Огородников~Е.~Н. Радченко~В.~П., Яшагин~Н.~С.
\paper Реологические модели вязкоупругого тела с~памятью и~дифференциальные уравнения дробных осцилляторов
\jour Вестн. Сам. гос. техн. ун-та. Сер. Физ.-мат. науки
\yr 2011
\issue 1(22)
\pages 255--268
\mathnet{http://mi.mathnet.ru/rus/vsgtu932}
\crossref{10.14498/vsgtu932}
\misctransl
\by Ogorodnikov~E.~N., Radchenko~V.~P., Yashagin~N.~S.
\paper Rheological model of viscoelastic body with memory and differential equations of fractional oscillator
\jour Vestn. Samar. Gos. Tekhn. Univ. Ser. Fiz.-Mat. Nauki \rm [J. Samara State Tech. Univ., Ser. Phys. \& Math. Sci.]
\yr 2011
\issue 1(22)
\pages 255--268
\lang In Russian


\RBibitem{ulg:bib:OY_some_properties_operators}
\by Огородников~Е.~Н., Яшагин~Н.~С.
\paper О некоторых свойствах операторов с~функциями типа Миттаг--Леффлера в ядрах
\inbook Труды шестой Всероссийской научной конференции с международным участием \rm (1--4 июня 2009~г.).Часть~3
\bookinfo Дифференциальные уравнения и краевые задачи
\serial Матем. моделирование и краев. задачи
\yr 2009
\pages 181--188
\publ СамГТУ
\publaddr Самара
\mathnet{http://mi.mathnet.ru/rus/mmkz1369}
\misctransl
\by Ogorodnikov~E.~N., Yashagin~N.~S.
\paper On some properties of operators with Mittag-Leffler type functions in kernels
\inbook Proceedings of the Sixth All-Russian Scientific Conference with international participation \rm (1--4 June 2009). Part~3
\serial Matem. Mod. Kraev. Zadachi
\yr 2009
\pages 181--188
\publ Samara State Technical Univ.
\publaddr Samara


\RBibitem{ulg:bib:Abusaitova_some_special_functions}
\by Абусаитова~Л.~Г., Огородников~Е.~Н.
\paper О некоторых специальных функциях, связанных с функцией Миттаг--Леффлера, их свойствах и применении
\inbook Нелокальные краевые задачи и проблемы современного анализа и информатики
\bookinfo Материалы X Школы молодых ученых
\yr 2012
\pages 13--15
\publ КБНЦ РАН
\publaddr Нальчик
\misctransl
\by Abusaitova~L.~G., Ogorodnikov~E.~N.
\paper Some special functions associated with the Mittag--Leffler function, their properties, and applications
\inbook Nelokal'nye kraevye zadachi i problemy sovremennogo analiza i informatiki \rm [Non-local boundary value problems and problems of modern analysis and informatics]
\yr 2012
\pages 13--15
\publaddr Nal'chik
\lang In Russian


\Bibitem{ulg:bib:Gorenflo-Mainardi1997}
\by Gorenflo~R., Mainardi~F.
\paper Fractional Calculus. Integral and Differential Equations of Fractional Order
\inbook Fractals and Fractional Calculus in Continuum Mechanics
\serial CISM Courses and Lectures
\yr 1997
\vol 378
\pages 223--276
\publ Springer
\publaddr Wien
\crossref{10.1007/978-3-7091-2664-6_5}
\zmath{https://zbmath.org/?q=an:01273494}

\Bibitem{ulg:bib:Koeller1984}
\by Koeller~R.~C.
\paper Applications of Fractional Calculus to the Theory of Viscoelasticity
\jour J.~Appl. Mech.
\yr 1984
\vol 51
\issue 2
\pages 299--307
\crossref{10.1115/1.3167616}

\Bibitem{ulg:bib:Carpinteri}
\by Carpinteri~A., Cornetti~P., Sapora~A.
\paper Nonlocal elasticity: an approach based on fractional calculus
\jour Meccanica
\yr 2014
\vol 49
\issue 11
\pages 2551--2569
\crossref{10.1007/s11012-014-0044-5}
\scopus{http://www.scopus.com/record/display.url?origin=inward&eid=2-s2.0-84911998390}

\Bibitem{ulg:bib:Bagley_theoretical_basis}
\by Bagley~R.~L., Torvic~P.~J.
\paper A Theoretical Basis for the Application of Fractional Calculus to Viscoelasticity
\jour J.~Rheol.
\yr 1983
\vol 27
\issue 3
\pages 201--210
\crossref{10.1122/1.549724}
\adsnasa{http://adsabs.harvard.edu/cgi-bin/bib_query?1983JRheo..27..201B}
\scopus{http://www.scopus.com/record/display.url?origin=inward&eid=2-s2.0-0020765202}

\Bibitem{ulg:bib:Bagley_Fractional_Calculus}
\by Bagley~R.~L., Torvic~P.~J.
\paper Fractional calculus --- A different approach to the analysis of viscoelastically damped structures
\jour AIAA Journal
\yr 1984
\vol 21
\issue 5
\pages 741--748
\crossref{10.2514/3.8142}
\adsnasa{http://adsabs.harvard.edu/cgi-bin/bib_query?1983AIAAJ..21..741B}

\Bibitem{ulg:bib:Lewandowski}
\by Lewandowski~R., \Russian Chor\k{a}\.zyczewski~B.
\paper Identification of the parameters of the Kelvin--Voigt and the Maxwell fractional models, used to modeling of viscoelastic dampers
\jour Computers and Structures
\yr 2009
\vol 88
\issue 1--2
\pages 1--17
\crossref{10.1016/j.compstruc.2009.09.001}
\scopus{http://www.scopus.com/record/display.url?origin=inward&eid=2-s2.0-70350770826}

\Bibitem{ulg:bib:CaputoMainardi_LinearModels}
\by Caputo~M., Mainardi~F.
\paper Linear models of dissipation in anelastic solids
\jour La Rivista del Nuovo Cimento
\yr 1971
\vol 1
\issue 2
\pages 161--198
\crossref{10.1007/bf02820620}
\adsnasa{http://adsabs.harvard.edu/cgi-bin/bib_query?1971NCimR...1..161C}
\scopus{http://www.scopus.com/record/display.url?origin=inward&eid=2-s2.0-51649182765}

\Bibitem{ulg:bib:CaputoMainardi_NewDissipation}
\by Caputo~M., Mainardi~F.
\paper A new dissipation model based on memory mechanism
\jour Pure and Applied Geophysics
\yr 1971
\vol 91
\issue 1
\pages 134--147
\crossref{10.1007/bf00879562}
\adsnasa{http://adsabs.harvard.edu/cgi-bin/bib_query?1971PApGe..91..134C}
\scopus{http://www.scopus.com/record/display.url?origin=inward&eid=2-s2.0-0000930143}

\Bibitem{ulg:bib:ScottBlair1947}
\by Scott Blair~G.~W.
\paper The role of psychophysics in rheology
\jour Journal of Colloid Science
\yr 1947
\vol 2
\issue 1
\pages 21--32
\crossref{10.1016/0095-8522(47)90007-x}

\Bibitem{ulg:bib:ScottBlair1949}
\by Scott Blair~G.~W.
\book A survey of general and applied rheology
\yr 1949
\publ Sir Isaac Pitman \& Sons, Ltd.
\publaddr London
\totalpages xvi+314
\zmath{https://zbmath.org/?q=an:0033.23202}

\RBibitem{ulg:bib:Gerasimov1948}
\by Герасимов А.~Н.
\paper Обобщение линейных законов деформирования и его применение к задачам внутреннего трения
\jour ПММ
\yr 1948
\vol 12
\issue 3
\pages 251--260
\zmath{https://zbmath.org/?q=an:0032.12901}
\misctransl
\by Gerasimov A.~N.
\paper A generalization of linear laws of deformation and its application to internal friction problem
\jour Prikl. Mat. Mekh.
\yr 1948
\vol 12
\issue 3
\pages 251--260
\lang In Russian


\Bibitem{ulg:bib:Barrett}
\by Barrett~J.~H.
\paper Differential equations of non-integer order
\jour Canad. J.~Math.
\yr 1954
\vol 6
\pages 529--541
\crossref{10.4153/cjm-1954-058-2}
\zmath{https://zbmath.org/?q=an:0058.10702}

\RBibitem{ulg:bib:AbusaitovaChuvashi2014}
\by Огородников~Е.~Н., Абусаитова~Л.~Г.
\paper Определяющие соотношения и~начальные задачи для вязкоупругих сред с~дробными операторами Римана--Лиувилля
\inbook Материалы VIII Всероссийской конференции по механике деформируемого тв\"eрдого тела
\yr 2014
\pages 105--107
\publ Чуваш. гос. пед. ун-т
\publaddr Чебоксары
\issueinfo Ч.~2
\misctransl
\by Ogorodnikov~E.~N., Abusaitova~L.~G.
\paper Relations definition, and initial problem for viscoelastic media with fractional Riemann--Liouville operators
\inbook The VIII All-Russian Conference on Mechanics Mechanics of Deformable Solids
\bookinfo Book of Abstracts and Conference Materials
\yr 2014
\pages 105--107
\publ Chuvash State Pedagogical Univ.
\publaddr Cheboksary
\lang In Russian
\issueinfo Part~2


\RBibitem{ulg:bib:AbusaitovaMatPhys2014}
\by Абусаитова~Л.~Г., Огородников~Е.~Н.
\paper Математическое моделирование вязкоупругих сред с памятью и задача параметрической идентификации дробных реологических моделей
\inbook Четвертая международная конференция «Математическая физика и ее приложения»
\bookinfo материалы конф.
\yr 2014
\pages 40--41
\eds чл.-корр.~РАН И.~В.~Волович; д.ф.-м.н., проф. В.~П.~Радченко
\publ СамГТУ
\publaddr Самара
\misctransl
\by Abusaitova~L.~G., Ogorodnikov~E.~N.
\paper Mathematical modeling of viscoelastic fluids with memory and the problem of parametric identification of fractional rheological models
\inbook The 4nd International Conference ``Mathematical Physics and its Applications''
\bookinfo Book of Abstracts and Conference Materials
\yr 2014
\pages 40--41
\eds I.~V.~Volovich; V.~P.~Radchenko
\publ Samara State Technical Univ.
\publaddr Samara
\lang In Russian


\RBibitem{ulg:bib:UngarovaKazan2015}
\by Унгарова~Л.~Г.
\paper Явные решения задачи о ползучести для некоторых нелинейных реологических моделей наследственно-упругого тела
\inbook ХI Всероссийский съезд по фундаментальным проблемам теоретической и прикладной механики
\bookinfo сборник докладов
\yr 2015
\pages 3843--3845
\publaddr Казань
\elib{http://elibrary.ru/item.asp?id=24825963}
\misctransl
\by Ungarova~L.~G.
\paper Explicit solutions of the creep problem for some non-linear rheological models of hereditary elastic body
\inbook The XI All-Russian Congress on Fundamental Problems of Theoretical and Applied Mechanics
\bookinfo Abstracts Book
\yr 2015
\pages 3843--3845
\publaddr Kazan'
\lang In Russian

\RBibitem{ulg:bib:RadchenkoGoludin}
\by Радченко~В.~П., Голудин~Е.~П.
\paper Феноменологическая стохастическая модель изотермической ползучести поливинилхлоридного пластиката
\jour Вестн. Сам. гос. техн. ун-та. Сер. Физ.-мат. науки
\yr 2008
\issue 1(16)
\pages 45--52
\crossref{10.14498/vsgtu571}
\misctransl
\by Radchenko~V.~P., Goludin~E.~P.
\paper Phenomenological stochastic isothermal creep model for an polivinylchloride elastron
\jour Vestn. Samar. Gos. Tekhn. Univ. Ser. Fiz.-Mat. Nauki \rm [J. Samara State Tech. Univ., Ser. Phys. \& Math. Sci.]
\yr 2008
\issue 1(16)
\pages 45--52
\lang In Russianxx+347


\end{thebibliography}


\newpage

\makeentitle

\vspace{-5mm}

\AdditionalContent{Competing interests}{Specify what conflicts of interest are available.} 
%\AdditionalContent{Competing interests}{I have ... (or) no competing interests.} 
%\AdditionalContent{Competing interests}{We have ... (or) no competing interests.} 
\AdditionalContent{Authors' contributions and responsibilities}{What is the author's responsibility for each author?}
%\AdditionalContent{Author's Responsibilities}{I take full responsibility for submitting the final manuscript in print. I approved the final version of the manuscript.} 
%\AdditionalContent{Authors' contributions and responsibilities}{Each author has participated in the article concept development and in the manuscript writing. The authors are absolutely responsible for submitting the final manuscript in print. Each author has approved the final version of manuscript.} 

\AdditionalContent{Funding}{Indicate the source(s) of your funding.}
%\AdditionalContent{Funding}{The research has not had any funding.}
%\AdditionalContent{Funding}{This work was supported by the Russian Foundation for Basic Research (project no. xx–xx–xxxxx_a).}

\AdditionalContent{Acknowledgments}{Express your gratitude (acknowledgments) here.}
%\AdditionalContent{Acknowledgments}{The authors are grateful to the  referee for  careful reading  of  the paper and valuable  suggestions and comments.}


\end{document}
